\chapter{Some Constructions}

% Try to weave references to the literature into the narrative
The problem of relating the structure of $\Gamma$ 
to that of the triple $(\Gamma_\cB, \Gamma[B,C], \cD(B))$
is very difficult when $\Gamma$ is taken to be 
an arbitrary imprimitive symmetric graph.
The approach which has been taken by researchers in the area
is to first impose additional assumptions 
on one or more of the components of the decomposition,
then study the subfamily of imprimitive symmetric graphs 
admitting a decompositions which satisfies these additional assumptions.

The special case where the parameters of the ``kernel design'' $\cD(B)$
satisfy: $$v = k-1 \geq 2.$$
was studied by Li, Praeger and Zhou \cite{ZCP2000}.
It was found that for this special case, 
if the design $\cD(B)$ contains no repeated blocks,
then there exists an elegant combinatorial method
for constructing $\Gamma$ from $\Gamma_\cB$.
This is the {\it three-arc graph construction} which will
be described in the next section.

In a later paper Zhou \cite{ZDM2002} showed that 
the three-arc graph construction actually applies 
to a wider family of triples $(G, \Gamma, \cB)$
than those originally studied.
In particular the construction may be used for any triple $(G, \Gamma, \cB)$
satisfying the following condition:  

\begin{condition}
The induced actions of $G_B$ on $B$ and $\Gamma_\cB(B)$
are permutationally equivalent with respect to some bijection
$\rho : B \to \Gamma_\cB(B)$.
\end{condition}

On a group theoretic level this is quite a natural condition to impose.
Suppose that the triple $(G,\Gamma, \cB)$ satisfies the above condition
and that $\Gamma \cong \Sab(G,K,a)$ 
and $\Gamma_\cB \cong \Sab(G,H,a)$.
If $B$ is a block of $\cB$ then
by proposition [?] the action of $G_B$ on $B$
is permutation equivalent to $(H, \cos(K))$.
The action of $G_B$ on $\Gamma_\cB(B)$ 
is just the ``local action'' of $\Gamma_\cB$
and so the permutation group $(G_B, \Gamma_\cB(B))$
is permutation equivalent to $(H, \cos(a^{-1}Ha \cap H))$.

If the condition above holds, 
then we must have $K \cong a^{-1}Ha \cap H$.
This means that the stabilizer of a point of $\Gamma$ is isomorphic
to the stabilizer of an arc of $\Gamma_\cB$.
This strong ``coupling'' between the points of $\Gamma$
and the arcs $\Gamma_\cB$ allows us to construct
$\Gamma$ from $\Gamma_\cB$ in a straightforeward manner.

%The condition above is also very natural on a design theoretic level.
%Recall that if $\Gamma \cong \Sab(G,K,a)$ and $\Gamma_\cB \cong \Sab(G,H,a)$
%then the ``kernel design'' of the homomorphism $\Gamma \mapsto \Gamma_\cB$
%has point stabilizer $a^{-1}Ha \cap H$ and block stabilizer $\cK$.
%If $\cK \cong a^{-1}Ha \cap H$ then the action of $H$ on the points
%of $\cK$ is permutation equivalent to the action of $H$ on the blocks
%of $\cK$.

\section{Three-Arc Graphs}
The three-arc graph construction was first introduced 
by Li, Praeger and Zhou in \cite{ZCP2000}.
Let $\Sigma$ be any $G$-symmetric graph.
Recall that an {\it $s$-arc} of $\Sigma$ is a sequence 
$(\alpha_0, \alpha_1, \dots, \alpha_s)$ of vertices in $\Sigma$ such that 
$\alpha_i,\alpha_{i+1}$ are adjacent in $\Sigma$ and $\alpha_{i-1} \neq \alpha_{i+1}$
for each $i$. The set of $s$-arcs of $\Sigma$ is denoted by $\Arc_s(\Sigma)$.
Consider the induced action of $G$ on $\Arc_3(\Sigma)$ given by:
$$(\sigma_0,\sigma_1,\sigma_2,\sigma_3)^g = (\sigma_0^g,\sigma_1^g,\sigma_2^g,\sigma_4^g).$$
This action is, in general, intransitive. We may however partition the set $\Arc_3(\Sigma)$
into {\it orbits} on which $G$ acts transitively. For such an orbit $\Delta$, let 
$\Delta^{\circ} = \{ (\sigma_4, \sigma_3, \sigma_2, \sigma_1) : 
                 (\sigma_1, \sigma_2, \sigma_3, \sigma_4) \in \Delta \}$ 
denote the {\it pair} of $\Delta$. It is not hard to check that
$\Delta^{\circ}$ is again an orbit of $G$ on $\Arc_3(\Sigma)$.
If $\Delta = \Delta^{\circ}$ then $\Delta$ is said to be {\it self paired}.

\begin{definition}
\label{threearc}
Given a $G$-symmetric graph $\Sigma$ and a self-paired orbit $\Delta$ on 
$\Arc_3(\Sigma)$, the {\it three-arc graph}
$\Gamma = \Arc_{\Delta}(\Sigma)$ is the graph with vertex set $A \Sigma$
and arc set $\{ ((\sigma,\tau), (\sigma',\tau') : (\tau,\sigma,\sigma',\tau') \in \Delta \}$.
\end{definition}

Note that the requirement that $\Delta$ is self-paired 
ensures that the resulting graph is simple.

\begin{proposition}
With $G$, $\Sigma$ and $\Delta$ as in definition \ref{threearc},
the three-arc graph $\Gamma = \Arc_{\Delta}(\Sigma)$
is $G$-symmetric.
\begin{proof} Immediate from the construction.
\end{proof} 
\end{proposition}

For each vertex $\sigma$ of $\Sigma$, let 
$B(\sigma) = \{ (\sigma, \tau) : \tau \mbox{ is a neighbour of } \sigma \}$
be the set of arcs of $\Sigma$ with initial vertex $\sigma$.
Clearly $\cB = \{ B(\sigma) : \sigma \mbox{ is a vertex of } \Sigma \}$ 
is a partition of $A \Sigma$.  
For any $g \in G$ we have 
\begin{eqnarray*}
B(\sigma)^g & = & \{ (\sigma, \tau)^g : \tau \mbox{ is a neighbour of } \sigma \} \\
              & = & \{ (\sigma^g, \tau^g) : \tau \mbox{ is a neighbour of } \sigma \} \\
	      & = & \{ (\sigma^g, \tau) : \tau \mbox{ is a neighbour of } \sigma^g \} \\
	      & = & B(\sigma^g).
\end{eqnarray*}
Thus $\cB = \{ B(\sigma) : \sigma \in V \Sigma \}$ 
is in fact a $G$-invariant partition of $V \Gamma$. 

\begin{proposition}
\label{three-arc quotient}
With $\Gamma = \Arc_{\Delta}(\Sigma)$ and $\cB$ as defined above,
The quotient graph $\Gamma_\cB$ is isomorphic to $\Sigma$.

\begin{proof}
The map $\sigma \mapsto B(\sigma)$ identifies the vertices 
of $\Sigma$ with the vertices of $\Gamma_\cB$.
We must show that 
$(\sigma, \sigma')$ is an arc of $\Sigma$ if and only if
$(B(\sigma), B(\sigma'))$ is an arc of $\Gamma_\cB$.

If $(B(\sigma), B(\sigma'))$ is an arc of $\Gamma_\cB$
then then for some $\tau, \tau' \in V \Sigma$ 
the arcs $(\sigma, \tau)$ and $(\sigma',\tau')$ are adjacent in $\Gamma$.
It follows that 
$(\tau, \sigma, \sigma',\tau')$ is a three-arc of $\Sigma$
and so $\sigma$ is adjacent to $\sigma'$ in $\Sigma$. 

For the other direction, 
let $(\alpha, \beta, \gamma, \delta)$ be any three-arc 
of $\Sigma$ contained in in $\Delta$.
If $(\sigma, \sigma')$ is any arc of $\Sigma$, 
then by the $G$-arc-transitivity of $\Sigma$
there exists a $g \in G$ such that $(\beta,\gamma)^g = (\sigma,\sigma')$.
Since $\Delta$ is an orbit of $G$ on the three-arcs of $\Sigma$, 
it follows that 
$(\alpha,\beta,\gamma,\delta)^g = (\alpha^g, \sigma, \sigma', \delta^g)$ 
is contained in $\Delta$.
Thus the arc $(\sigma, \alpha^g)$ is adjacent to 
the arc $(\sigma', \delta^g)$ in $\Gamma$
and so $B(\sigma)$ is adjacent to $B(\sigma')$ in $\Gamma_\cB$.


\end{proof}
\end{proposition}

%The three-arc graph construction allows us to take a $G$-symmetric graph
%$\Sigma$ and construct a {\it larger} imprimitive $G$-symmetric graph $\Gamma$
%which admits the original graph $\Sigma$ as a quotient.
%Provided of course that we can find a suitable self-paired orbit
%on the three arcs of $\Sigma$.
%This is an interesting result because blah blah blah....

%We shall now show that the family of triples $(G, \Gamma, \cB)$
%for which $\Gamma$ is a three-arc graph of $\Gamma_\cB$ 
%is precisely the family of triples satisfying 
%the condition {\bf (PE)} stated earlier.



\begin{proposition}

If $(G,\Gamma, \cB)$ is such that $\Gamma$ is a three-arc graph
of $\Gamma_\cB$, then for any $B \in \cB$
the permutation groups $(G, B)$ and $(G, \Gamma_\cB(B))$
are permutation equivalent.  

\begin{proof}
Each $B \in \cB$ is of the form $B(\sigma)$ for some
$\sigma$ a vertex of $\Gamma_\cB$.
The vertices of $\Gamma$ contained in $B(\sigma)$ 
are precisely the arcs of $\Gamma_\cB$
of the form $(\sigma, \tau)$.

Let $\rho : B(\sigma) \to \Gamma_\cB(B(\sigma))$ be given by 
$(\sigma,\tau) \mapsto B(\tau)$.
Clearly $\eta$ is a bijection.
If $g$ is any element of the stabilizer of $B(\sigma)$ 
then we have:
\begin{eqnarray*}
\rho((\sigma,\tau)^g) & = & \rho((\sigma^g,\tau^g)) \\
                      & = & \rho((\sigma,\tau^g))   \\
		      & = & B(\tau^g)               \\
		      & = & B(\tau)^g.
\end{eqnarray*}
Thus $\rho$ induces a permutation equivalence between the permutation groups
$(G_{B(\sigma)}, B(\sigma))$ and $(G_{B(\sigma)}, \Gamma_\cB(B(\sigma))$.

\end{proof}
\end{proposition}

\begin{proposition}
If $(G, \Gamma, \cB)$ is such that 
$(G, B)$ is permutation equivalent to $(G, \Gamma_\cB(B))$
for some $B \in \cB$, then the action of $G$
on the vertices of $\Gamma$ is permutation equivalent
to the action of $G$ on the arcs of $\Gamma_\cB$.

\begin{proof}
Let $\rho : B \to \Gamma_\cB(B)$ be any bijection
inducing a permutation equivalence between 
$(G,B)$ and $(G, \Gamma_\cB(B))$
and let $\pi : V \Gamma \to \cB$ be the map which sends
each vertex of $\Gamma$ to the block of $\cB$ containing it.
Fix a vertex $\alpha$ of $\Gamma$.
Since $\Gamma$ is $G$-vertex transitive, every vertex of $\Gamma$
may be written in the form $\alpha^g$ for some $g \in G$.
Let $\mu : V \Gamma \to A \Gamma_{\cB}$ be given by
$ \alpha^g \mapsto (\pi(\alpha)^g, \rho(\alpha)^g) $.
It is not too hard to see that $\mu$ is a bijection.
For any $\beta = \alpha^g$ and any $x \in G$ 
we have:
\begin{eqnarray*}
\mu(\beta^x) & = & \mu(\alpha^{gx})                       \\
             & = & (\pi(\alpha)^{gx}, \eta(\alpha)^{gx})    \\
	     & = & (\pi(\alpha^g)^x, \eta(\alpha^g)^x)      \\
	     & = & (\pi(\beta)^x, \eta(\beta)^x)            \\
	     & = & (\pi(\beta), \eta(\beta))^x              \\
	     & = & \mu(\beta)^x
\end{eqnarray*}
Thus $\mu$ establizhes a permutation isomorphism between
$(G,V \Gamma)$ and $(G, A \Gamma_\cB)$.

\end{proof}
\end{proposition}

We may make use the map $\mu$ to ``label'' the vertices
of $\Gamma$ by the arcs of $\Gamma_\cB$.
For any arc $(B,C)$ of $\Gamma_{\cB}$, 
let $v_{BC} = \mu^{-1}((B,C))$ denote
the vertex of $\Gamma$ which is mapped to $(B,C)$ by $\mu$.

\begin{proposition}
\label{labelling}
Some stuff about how $G$ acts on labelled vertices
and how the initial block contains the vertex.
\end{proposition}

\begin{proposition}
\label{PE implies 3-arc}
Provided that $\Gamma$ as valency at least two,
for each arc $(v_{BC}, v_{DE})$ of $\Gamma$,
$(C,B,D,E)$ is a 3-arc of $\Gamma_\cB$.

\begin{proof}
Suppose that $v_{BC}$ were adjacent to $v_{CB}$.
Since $\val(\Gamma) \geq 2$, the vertex $v_{CB}$
is adjacent to some other vertex $v_{B_1C_1}$ 
distinct from $v_{BC}$
By the $G$-symmetry of $\Gamma$ there exists a $g \in G$
such that $(v_{BC},v_{CB})^g = (v_{B_1C_1},v_{CB})$.
By the previous proposition, this implies that 
$B = B^x = B_1$ and $C = C^x = C_1$, so $v_{BC} = v_{B_1C_1}$.
A contradiction. Thus $v_{BC}$ is not adjacent to $v_{CB}$.

Suppose now that $v_{BC}$ were adjacent to $v_{CE}$
for some $E \neq B$.
Since $\val(\Gamma) \geq 2$ we can find another vertex $v_{C_1E_1}$
distinct from $v_{CE}$ and adjacent to $v_{BC}$.
By the $G$-symmetry of $\Gamma$ we can find some $g_1 \in G$
such that $(v_{BC},v_{CE})^{g_1} = (v_{BC},v_{C_1E_1})$.
By the previous proposition, this implies that $C = C^{g_1} = C_1$.
We can also find a $g_2 \in G$
such that $(v_{BC},v_{CE})^{g_2} = (v_{C_1E_1}, v_{BC})$
and so 
$B = C^{g_2} = E_1$. Combining these gives $v_{C_1E_1} = v_{CB}$, 
but $v_{BC}$ is not adjacent to $v_{CB}$
so again we have a contradiction.

We know now that if $(v_{BC}, v_{DE})$ is an arc of $\Gamma$,
then $D \neq C$. A simmilar argument shows that $B \neq E$,
so $B,C,D,E$ are distinct vertices with $B$ adjacent to $C$
and $D$ adjacent to $E$. Since, by the previous proposition
$B(v_{BC}) = B$ and $B(v_{DE} = D$, we have $B$ adjacent to $D$
and so $C,B,D,E$ is a 3-arc of $\Gamma_\cB$ as claimed.

\end{proof}
\end{proposition}

Concluding remarks.

%It follows immediately from proposition \ref{PE implies 3-arc}
%that if $(\Gamma, \cF)$ is any pair satisfying {\bf PE assumption}
%then $\Gamma$ is isomorphic to a 3-arc graph of $\Gamma_\cF$
%taking $\Delta$ to be the orbit containing $(C,B,D,E)$
%for any arc $(v_{BC}, v_{DE})$ of $\Gamma$.
%(This needs to be much better written).
%Need to talk about conditions under which a $G$-symmetric graph
%admits a self-paired $G$-orbit on its three-arcs.

%\section{$(G,2)$-arc transitive quotients}

%\section{$(G,3)$-arc transitive quotients}

%\section{Almost Covers}

%\section{The case when the ``kernel design'' is symmetric}

%\section{Uniqueness of the Construction}

%\begin{enumerate}
%\item how to find the self-paired orbit?
%\item not unique
%\item extremal case: induced bipartite graph almost a matching
%\item occurs only if quotient is $(G,3)$-arc transitive
%\item which paper does this occur in?
%\end{enumerate}

%We have seen that any pair $(\Gamma, \cB)$ satisfying the
%{\it (PE) assumption}, $\Gamma$ may be reconstructed
%from $\Gamma_\cB$ by using the three-arc graph construction
%{\it for some suitable self-paired orbit}. 
%The question is, which orbit? and is there a unique such orbit?

%An obvious case in which there will be a unique orbit
%on the three arcs is when the quotient is in fact
%$(G,3)$-arc transitive. In this case, $G$ is in fact
%transitive on $\Arc_3(\Gamma_\cB)$ and the set
%of all $3$-arcs forms a self-paired orbit. 
\section{Locally imprimitive quotient}

Suppose that $\Gamma = \sab(G,H,a)$.
Let $\overline{H} = a^{-1}Ha \cap H$,
so the local permutation group induced at each vertex of $\Gamma$
is equivalent to $(H, \cos(\overline{H}))$.
Suppose further that there is some $K$ such that $\overline{H} < K < H$
and $a \not\in K$. Then the graph $\widetilde{\Gamma} = \sab(G,K,a)$
is an extension of $\Gamma$.
Let $\overline{K} = a^{-1}Ka \cap K$.
Since $K < H$ we must also have $a^{-1}Ka < a^{-1}Ha$
so $\overline{K} \leq \overline{H}$.
Since $\overline{H} < K$ we must also have
$a^{-1}\overline{H}a < a^{-1}Ka$
but since $a$ is an involution $a^{-1}\overline{H}a = \overline{H}$
so we have $\overline{H} \leq \overline{K}$.
That is $\overline{H} = \overline{K}$
and the local permutation group induced at each vertex of
$\widetilde{\Gamma}$ is $(K, \cos(\overline{H}))$.


The pair $(\Gamma,\widetilde{\Gamma})$ satisfy the property that
{\it globally} $\widetilde{\Gamma}$ is imprimitive,
admitting $\Gamma$ as quotient, but {\it locally}
$\Gamma$ is imprimitive, and the local permutation group of
$\widetilde{\Gamma}$ is a quotient of the local permutation
group of $\Gamma$. We wish to understand how $\Gamma$
and $\widetilde{\Gamma}$ are related {\it structurally},
and ideally find some method of constructing $\widetilde{\Gamma}$
from $\Gamma$.

As a preliminary observation, let $n = [G:H]$ be the number of vertices of $\Gamma$,
let $v = [K:\overline{H}]$ be the valency of $\widetilde{\Gamma}$
and let $r = [H:K]$.
The graph $\widetilde{\Gamma}$ has $[G:K] = [G:H][H:K] = nr$
vertices. That is, $r$ times as many vertices as $\Gamma$.
The valency of $\Gamma$ is $[H:\overline{H}] = [H:K][K:\overline{H}] = rv$.
That is $r$ times the valency of $\widetilde{\Gamma}$.
Since the number of edges in a graph is equal to half the
valency times the number of vertices, it follows that
both $\Gamma$ and $\widetilde{\Gamma}$ have the same number of edges.
\hfill \break


The action of $G$ on the arcs of $\Gamma$ is permutation equivalent
to $(G, \cos(\overline{H}))$.
Since $\Gamma$ is simple, each arc $(\alpha,\beta)$ has a {\it pair}
namely $(\beta,\alpha)$.
The function $\varphi : A \Gamma \to A \Gamma$
which sends each arc to its pair is an involution
It also preserves the action of $G$.

Now, since $\overline{H} < H < G$ the action of $G$
on the arcs of $\Gamma$ is imprimitive.
That is, there is some $G$-invariant partition $\cB$
of $A \Gamma$ such that the action of $G$ on $\cB$
is permutation equivalent to $(G, \cos(H))$.
In fact, this is the partition $\cB = \{ B(\alpha) : \alpha \in V \Gamma \}$
where $B(\alpha)$ is the set of all arcs who's initial vertex is $\alpha$.

Let $\pi : A \Gamma \to \cB$ be the map which sends
each arc to the block of $\cB$ containing it.
We may define a new graph $\Sigma$ we vertex set
$\cB$ where $p$ is adjacent to $b$ if and only if
there is some $x \in \pi^{-1}(p)$ and some $y \in \pi^{-1}(b)$
such that $y = \varphi(x)$.
This new graph $\Sigma$ is in fact isomorphic to $\Gamma$.


The action of $G$ on the arcs of $\Gamma$ is permutation
isomorphic to the action of $G$ on the arcs of $\widetilde{\Gamma}$.
Both actions are of the form $(G, \cos(\overline{H}))$
Since $\overline{H} < K < G$ there exists some $G$-invariant partition
$\mathfrak{B}$ of $A \Gamma$ such that the action of $G$ on $\mathfrak{B}$
is permutation equivalent to $(G, \cos(H))$.
The partition $\mathfrak{B}$ is a refinement of the partition $\cB$.

Let $\lambda : A \Gamma \to \mathfrak{B}$ be the map hich sends
each arc to the block of $\mathfrak{B}$ containing it.
We may define a new graph $\widetilde{\Sigma}$ we vertex set
$\mathfrak{B}$ where $p$ is adjacent to $b$ if and only if
there is some $x \in \lambda^{-1}(p)$ and some $y \in \lambda^{-1}(b)$
such that $y = \varphi(x)$.
This new graph $\Sigma$ is in fact isomorphic to $\widetilde{\Gamma}$.
Thus we have a method of construction $\widetilde{\Gamma}$ from $\Gamma$.

\section{Covering Graphs}

% Introductory paragraph:
The three-arc graph construction allows us, under certain conditions,
to ``unfold'' a $G$-symmetric graph into a larger, {\it imprimitive}
$G$-symmetric graph, admitting the original graph as a quotient.
The {\it covering graph} construction \cite{B1974} is simmilar
in this respect. It also, under certain conditions,
allows us to ``unfold'' a $G$-symmetric graph into a larger one.
It differs from the three-arc graph construction, however,
in that it requires a simultaneous ``unfolding''
of the group $G$.


% Semidirect Product:
Recall that given two groups $N$ and $G$ and a homomorphism
$$\rho : G \to \Aut(N)$$
we may form the {\it semidirect product} of $N$ by $G$.
This is the group $\widetilde{G} = N \rtimes_\rho G$,
whose elements are the ordered pairs:
$$\{ (n,g) : n \in N, g \in G \}$$
with multiplication given by:
$$(n_1, g_1)(n_2,g_2) = (n_1^{\rho(g_1)}n_2, g_1g_2).$$
The functions $i_1 : N \to \widetilde{G}$ and $i_2 : G \to \widetilde{G}$
given by $ n \mapsto (n,1) $ and $ g \mapsto (1,g) $ respectively give
natural embeddings of $N$ and $G$ into $\widetilde{G}$.
Identifying $N$ with its image under $i_1$,
the semidirect product has the property that $N$ is normal in $\widetilde{G}$
and $\widetilde{G}/N \cong G$.

% Definitions:
\begin{definition} [$N$-chain]
Suppose that $\Gamma$ is a $G$-symmetric graph and $N$ is a group.
An $N$-chain is a function $\phi : A \Gamma \to N$
satisfying $\phi((v,u)) = \phi((u,v))^{-1}$ for all $(u,v) \in A \Gamma$.
\end{definition}

\begin{definition} [Compatible $N$-chain]
Suppose that $\Gamma$ is a $G$-symmetric graph, $N$ is a group
and $\rho : G \to \Aut(N)$ is a homomorphism.
An $N$-chain $\phi$ is said to be {\it compatible} with $\rho$
if for every $g \in G$ the following diagram commutes:

\centerline{
\xymatrix{
\ar[d]_{g} \ar[r]^\phi A \Sigma & K \ar[d]^{\rho(g)} \\
\ar[r]_{\phi} A \Sigma & K
}}
\end{definition}

\begin{definition}[Biggs Cover]
Suppose that $\Gamma$ is a $G$-symmetric graph, $N$ is a group,
$\rho : G \to \Aut(N)$ is a homomorphism and $\phi$ is a compatible $N$-chain.
The Biggs Cover $\widetilde{\Gamma}(N,\rho,\phi)$
of $\Gamma$ with respect to $\phi$ is the graph
with vertex set:
$\{ (g,v) : k \in G, v \in V \Gamma \}$
and arc set:
$ \{ ((g_1,v_1),(g_2,v_2)) : (v_1,v_2) \in A \Gamma, g_2 = g_1 \phi((v_1,v_2)) \}$
\end{definition}

Note that the condition $\phi((v,u)) = \phi((u,v))^{-1}$ ensures that the resulting
graph is simple.

% Proposition:
\begin{proposition}
The covering graph $\widetilde{\Gamma}(N,\rho,\phi)$ is $(N \rtimes_\rho G)$-symmetric.
\begin{proof}

Define an action of $N \rtimes_\rho G$ on $\widetilde{\Gamma}$ by:
$$(n,v)^{(\eta, g)} = (n^{\rho(g)} \eta,v^g).$$
This action is well-defined since:
\begin{eqnarray*}
(n,v)^{(\eta_1, g_1)(\eta_2, g_2)}
& = & (n^{\rho(g_1)} \eta_1 ,v^{g_1})^{(\eta_2, g_2)} \\
& = & (n^{\rho(g_1) \rho(g_2)} \eta_1^{\rho(g_2)} \eta_2 , v^{g_1 g_2}) \\
& = & (n,v)^{(\eta_1^{\rho(g_2)} \eta_2, g_1 g_2)}
\end{eqnarray*}
We must check that it preserves the adjacency structure of $\widetilde{\Gamma}$.
Observe that, by the compatibility of $\rho$ and $\phi$,
if $n_2 = n_1 \phi(v_1,v_2)$ then:
\begin{eqnarray*}
{n_2}^{\rho(g)} \phi({v_1}^g, {v_2}^g)
& = & {n_2}^{\rho(g)} \phi(v_1, v_2)^{\phi(g)} \\
& = & (n_2 \phi(v_1, v_2))^{\phi(g)} \\
& = & {n_1}^{\phi(g)}
\end{eqnarray*}
for any $(\eta,g) \in N \rtimes_\rho G$.
Also, since $G$ is a group of automorphisms of $\Gamma$, we know that
$(v_1,v_2) \in A \Gamma$ if and only if $({v_1}^g,{v_2}^g) \in A \Gamma$.
Thus the action defined above does indeed preserves adjacency.

\hfill \break
For any $(n_1,v_1), (n_2,v_2) \in V \widetilde{\Gamma}$,
by the $G$-symmetry of $\Gamma$ we can find a $g \in G$
such that ${v_1}^g = {v_2}$.
Let $\eta = {n_1}^{\rho(g)^{-1}} n_2$.
We have:
\begin{eqnarray*}
(n_1,v_1)^{(\eta,g)} & = & ({n_1}^{\rho(g)} \eta,{v_1}^g) \\
                    & = & ( {n_1}^{\rho(g)}, {n_1}^{\rho(g)^{-1}} n_2,v_2) \\
		    & = & (n_2,v_2).
\end{eqnarray*}
Thus the action is vertex transitive.

\hfill \break
Suppose that $(n_1,v_1)$ and $(n_2,v_2)$ are both adjacent to $(n,g)$,
so $n_1 = n \phi(v,v_1)$ and $n_2 = n \phi(v,v_2)$.
By the $G$-symmetry of $\Gamma$ we can find a $g \in G_v$ such that
$v_1^g = v_2$.
Let $\eta = (n^{-1})^{\rho(g)}n$. We have:
\begin{eqnarray*}
(n,v)^{(\eta,g)} & = & (n^{\rho(g)} \eta, v^g) \\
                 & = & (n^{\rho(g)} (n^{-1})^{\rho(g)}n, v) \\
		 & = & (n,v)
\end{eqnarray*}
and:
\begin{eqnarray*}
(n_1,v_1)^{(\eta,g)} & = & ({n_1}^{\rho(g)} \eta, {v_1}^g) \\
                 & = & ({n_1}^{\rho(g)} (n^{-1})^{\rho(g)}n, v_2) \\
                  & = & ({(n_1 n^{-1})}^{\rho(g)} n, v_2) \\
                   & = & (\phi(v,v_1)^{\rho(g)} n, v_2) \\
                    & = & (\phi(v,v_2) n, v_2) \\
                     & = & (n_2 n^{-1} n, v_2) \\
                      & = & (n_2, v_2) \\
\end{eqnarray*}
Thus the action is locally transitive. The result follows.	      
\end{proof}							
\end{proposition}
												 In fact, a stronger result than this holds.
												 If $\Gamma$ is $(G,s)$-arc transitive
												 then $\widetilde{\Gamma}$ is $(\widetilde{G},s)$-arc transitive,
												 where $\widetilde{G} = N \rtimes_\rho G$.
												 The proof is by induction and may be found in \cite{B1974}.
												 This construction was originally used by Conway
												 to produce an infinitely family of $5$-arc transitive graphs.

												 \hfill \break
												 For each $v \in V \Gamma$, let $B(v) = \{ (n,v) : n \in N \}$.
												 For any $(\eta, g) \in \widetilde{G}$ we have:
												 \begin{eqnarray*}
												 B(v)^{(\eta,g)} & = & \{ (n,v)^{(\eta,g)} : n \in N \} \\
												                 & = & \{ (\eta n^{\rho(g)},v^g) : n \in N \} \\
														                 & = & \{ (n, v^g) : n \in N \} \\
																                 & = & B(v^g).
																		 \end{eqnarray*}
																		 Thus the partition:
																		 $\cB = \{ B(v) : v \in V \Gamma \}$
																		 is $\widetilde{G}$-invariant.

\begin{proposition}
The quotient of the Biggs cover $\widetilde{\Gamma}$
with respect to the partition $\cB$
is isomorphic to the original graph $\Gamma$
\begin{proof}
The map $\eta : V \Gamma \to V \widetilde{\Gamma}_\cB$ given by
$v \mapsto B(v)$ establishes a bijection between the vertices of $\Gamma$
and the vertices of $\widetilde{\Gamma}_\cB$.
If $(u,v)$ is an arc of $\Gamma$, then $((1,u),(\phi((u,v)),v))$
is an arc of $\widetilde{\Gamma}$, and so $(B(u), B(v))$
is an arc of $\widetilde{\Gamma}_\cB$.
If $(u,v)$ is not an arc of $\Gamma$, then for all $n_1, n_2 \in N$
$((n_1, u),(n_2, v))$ is not an arc of $\widetilde{\Gamma}$
and so $(B(u), B(v))$ is not an arc of $\widetilde{\Gamma}_\cB$.
Thus $\eta$ establishes an isomorphism between
$\Gamma$ and $\widetilde{\Gamma}_\cB$.

\end{proof}
\end{proposition}
Let $(B,C)$ be any arc of $\widetilde{\Gamma}$.
By the above $B = B(u)$ and $C = B(v)$ for some $(u,v) \in A \Gamma$.
Its not too hard to see that each $(n,u) \in B$
has a unique neighbour in $C$, namely $(n \phi(u,v),v)$.
Thus the induced bipartite graph $\Gamma[B,C]$ is a matching.
It follows immediately that the valency of $\widetilde{\Gamma}$
is the same as the valency of $\Gamma$.

\section{Group Theoretic analysis of covering graph construction}

%The group $\widetilde{G}$ acts unfaithfully on $\cB$ with kernel
%$$\widetilde{N} = \{ (n,1) : n \in N \} \cong N.$$
%If $|N| = n$ then $|\widetilde{G}| = n |G|$,
%and since $\widetilde{\Gamma}$ has exactly $n$ times as many vertices as $\Gamma$,
%we have $[\widetilde{G} : \widetilde{G}_{B(\alpha)}] = n [G : G_\alpha]$.
%It follows that
%$ |\widetilde{G}_{B(\alpha)}| = |G_\alpha|$.

%Let $\pi : \widetilde{G} \to G$ be the natural projection.
%Clearly $|\pi(\widetilde{G}_{B(\alpha)})| \leq |\widetilde{G}_{B(\alpha)}|$.
%Since $G_\alpha \leq \pi(\widetilde{G}_{B(\alpha)})$,
%we must have $G_\alpha = \pi(\widetilde{G}_{B(\alpha)})$
%and $|\pi(\widetilde{G}_{B(\alpha)})| = |\widetilde{G}_{B(\alpha)}|$,
%so $\widetilde{G}_{B(\alpha)} \cong \pi(\widetilde{G}_{B(\alpha)}) \cong G_\alpha$.
%In fact, the permutation groups
%$(G_\alpha, \Gamma(\alpha))$ and
%$(\widetilde{G}_{B(\alpha)}, \widetilde{\Gamma}(B(\alpha)))$
%are equivalent via the bijection
%$\eta : \Gamma(\alpha) \to \widetilde{\Gamma}(B(\alpha))$
%which sends $\beta$ to $B(\beta)$.
%Probably a much easier way to say this.

%If $\Gamma$ is isomorphic to the coset graph $\Gamma(G, H, a)$
%then since the local actions of $\Gamma$ and $\widetilde{\Gamma}$
%are the same, $\widetilde{\Gamma}$ must be isomorphic to the coset graph
%$\Gamma(\widetilde{G}, H, \widetilde{a})$ where $\widetilde{a}$
%is such that $\pi(\widetilde{a}) = a$. That is
%$\widetilde{a} \in \{ (n,a) : n \in N \}$.
%Thus for any $N$ there are exactly $n$ possible Biggs covers of $\Gamma$.

%Looking closely at the compatibility condition between
%$\rho$ and $\phi$ we see that, since $G$ acts transitively
%on the arcs of $\Gamma$, the map $\phi$ is completely
%determined by where it sends the arc $(v_{H}, v_{Ha})$.
%Choosing where to send this arcs is essentially equivalent to
%choosing which preimage of $a$ we want.
%Requiring that $\phi((v_{H}, v_{Ha})) = n$ is essentially
%the same as determining that $\widetilde{a} = (n,a)$.

The Bigg's covering graph construction encompasses
all pairs of graphs $\Gamma(G, H, a)$ and $\Gamma(\widetilde{G}, H, \widetilde{a})$
where $\widetilde{G}$ is a semidirect product of $N$ by $G$ for some $N$.
The local permutation groups of $\Gamma$ and $\widetilde{\Gamma}$
are the same, and $N$ acts regularly on the fibers.

%In the more general case where $\widetilde{G}$ is an extention of $N$ by $G$
%but not a split extension $\Gamma(\widetilde{G}, H, \widetilde{a})$
%should be a cover of $\Gamma(G, H, a)$ in the sense that the induced
%bipartite graph is a matching, but Bigg's covering graph construction
%cannot be used to construct  $\Gamma(\widetilde{G}, H, \widetilde{a})$
%from $\Gamma(G, H, a)$.

%It should also possible to have a graph $\widetilde{\Gamma}$
%covering $\Gamma$ such that the local actions are the same degree
%but not permutation equivalent.

\section{Subgraph Extension}

I have been thinking about the idea of a ``combinatorial''
description of the extension of a graph in terms of its quotient.
That is, a description of the adjacency structure of $\widetilde{\Gamma}$
directly in terms of the adjaceny structure of $\Gamma$.
This section is a sketch of an idea I had, that perhaps
there is a connection between the subgraph structures
of different $G$-symmetric graphs for the same $G$.

\begin{definition} [Subgraph]
Suppose that $\Gamma = (V,A)$ is a graph.
If $W$ is a subset of $V$
and $B$ is a subset of $(W \times W) \cap A$
then $\Upsilon = (W,B)$ is a {\it subgraph} of $\Gamma$.
We write $\Upsilon < \Gamma$ and allow for the possibility
that $\Upsilon$ is a directed graph.
\end{definition}


\begin{definition} [Stabilizer of a Subgraph]
Suppose that $(G, \Gamma)$ is a symmetric graph and
$\Upsilon$ is a directed subgraph. The {\it stabilizer}
of $\Upsilon$ in $(G, \Gamma)$ is $\Aut(\Upsilon) \cap G$.
It is denoted by $G_{\Upsilon}$.
\end{definition}


\begin{definition} [Subgraph Graph]
Suppose that $(G, \Gamma)$ is a symmetric graph,
$\Upsilon$ is a directed subgraph of $\Gamma$ 
and $a$ is an involution of $\Gamma$ which fixes an arc.
The subgraph graph  $\sub(\Gamma, \Upsilon, a)$
of $\Gamma$ with respect to $\Upsilon$ and $a$
is the graph with vertices  $V = \{ \Upsilon^g : g \in G \}$
and arcs $A = \{ (\Upsilon \cup \Upsilon^a)^g : g \in G \}$.
\end{definition}

\begin{proposition}
Subgraph graph's are symmetric, with point stabilizer $G_{\Upsilon}$.
\begin{proof}
Obvious.
\end{proof}
\end{proposition}

\begin{example} [cube from tetrahedron]
Let $\Gamma$ be the tetrahedron with $S4$ acting on it.
Let $\Upsilon$ be any directed three-cycle.
Let $a$ be any involution fixing an arc of $\Upsilon$,
then the subgraph graph $\sub(\Gamma,\Upsilon,a)$
is isomorphic to the cube.
\end{example}


\begin{definition} [Symmetric subgraph]
Suppose that $(G,\Gamma)$ is a symmetric graph.
The digraph $\Upsilon < \Gamma$ is a {\it symmetric subgraph}
of $\Gamma$ if $(G_\Upsilon, \Upsilon)$ is a symmetric graph.
\end{definition}

\begin{question}
If $(G, \Gamma)$ is a symmetric graph, for which subgroups $K$ of $G$
do their exist a subgraphs of $\Upsilon$ of $\Gamma$ such that
$G_\Upsilon = K$ ?
For which subgroups $K$ of $G$
do their exist a {\it symmetric subgraph} $\Upsilon$ of $\Gamma$ such that
$G_\Upsilon = K$ ?
\end{question}

\begin{question}
Which extensions of a given $G$-symmetric graph $\Gamma$
are isomorphic to subgraph graphs of $\Gamma$ ?
Which are isomorphic to subgraph graphs of $\Gamma$
with respect to some symmetric subgraph (digraph) of $\Gamma$ ?
\end{question}



